\section*{الفصل \ref{ch:g12:space} --- المتجهات والمستقيمات والمستويات في الفضاء}

\begin{solution}{exo:g12:space:1}
$\vect{AB}(2, 1, -1)$ و $\vect{AC}(1, -1, 1)$ ليستا متناسبتين
($\frac21 \neq \frac{1}{-1}$)، إذن النقط ليست على استقامة واحدة. والمستقيم $(AB)$:
\[
(x, y, z) = (1 + 2t,\; t,\; 2 - t), \qquad t \in \R .
\]
\end{solution}

\begin{solution}{exo:g12:space:2}
$2(x - 1) - (y - 1) + 3(z - 0) = 0$، \emph{أي}
\[
2x - y + 3z - 1 = 0 .
\]
ومن أجل $B(0,2,1)$: $0 - 2 + 3 - 1 = 0$: نعم، تقع $B$ على المستوي.
\end{solution}

\begin{solution}{exo:g12:space:3}
$\vect{AB}(1,1,0)$ و $\vect{AC}(1,0,1)$:
\[
\cos\widehat{A}
= \frac{\vect{AB}\cdot\vect{AC}}{\norm{\vect{AB}}\,\norm{\vect{AC}}}
= \frac{1}{\sqrt2\,\sqrt2} = \frac12,
\]
إذن $\widehat A = 60^\circ$.
\end{solution}

\begin{solution}{exo:g12:space:4}
نقط $\mathcal D$: $(1 + t,\, 2 - t,\, 2t)$. وبالتعويض في
\omterm{def:g10:algebra:equation}{معادلة} $\mathcal P$:
\[
2(1+t) + (2 - t) - 2t + 1 = 5 - t = 0 \iff t = 5 .
\]
فنقطة \omterm{def:g10:numbers:interunion}{التقاطع} الوحيدة: $(6, -3, 10)$.
\end{solution}

\begin{solution}{exo:g12:space:5}
المتجهتان الناظمتان $(1, -1, 2)$ و $(2, 1, 1)$ ليستا \omterm{def:g12:space:collinear}{مستقيميتين}، إذن
يتقاطع المستويان في مستقيم. وبجمع المعادلتين: $3x + 3z = 5$، إذن
$x = \frac53 - z$؛ ثم تعطي \omterm{def:g10:algebra:equation}{المعادلة} الأولى
$y = x + 2z - 1 = \frac23 + z$. ومع $z = t$:
\[
(x, y, z) = \left(\frac53 - t,\; \frac23 + t,\; t\right), \qquad t \in \R,
\]
وهو مستقيم \omterm{def:g11:vect:vector}{متجهة} توجيهه $(-1, 1, 1)$. (والمعادلتان تتحققان
تطابقيًا.)
\end{solution}

\begin{solution}{exo:g12:space:6}
التوجيهان $\vec u_1(1, 2, -1)$ و $\vec u_2(1, 1, 2)$ ليسا مستقيميين،
إذن المستقيمان ليسا متوازيين. وأما \omterm{def:g10:numbers:interunion}{التقاطع} فيقتضي
\[
1 + t = 2 + s, \qquad 2t = 1 + s, \qquad 3 - t = 1 + 2s .
\]
فتعطي المعادلتان الأوليان $t = 1 + s$ و $2(1+s) = 1 + s$، إذن $s = -1$ و $t = 0$؛
لكن الثالثة تصير عندئذٍ $3 = -1$، وهذا سخيف. فلا نقطة مشتركة ولا
توازٍ: إذن المستقيمان متخالفان.
\end{solution}

\begin{solution}{exo:g12:space:7}
\emph{1.} $\operatorname{dist} = \dfrac{\abs{2\cdot3 - 1 + 2\cdot1 - 3}}
{\sqrt{4 + 1 + 4}} = \dfrac{4}{3}$.

\emph{2.} $H = M_0 + t\,\vec n$ مع $\vec n(2, -1, 2)$، باختيار $t$ بحيث
$H \in \mathcal P$:
\[
2(3 + 2t) - (1 - t) + 2(1 + 2t) - 3 = 4 + 9t = 0 \iff t = -\frac49 .
\]
ومنه $H\left(\frac{19}{9}, \frac{13}{9}, \frac19\right)$، و
\[
M_0H = \norm{t\,\vec n} = \frac49 \times 3 = \frac43,
\]
وهذا يوافق صيغة المسافة.
\end{solution}

\begin{solution}{exo:g12:space:8}
\omterm{def:g10:coordgeom:system}{الإحداثيات}: $G(1,1,1)$، والمستوي $(BDE)$ يمر بالنقط $B(1,0,0)$
و $D(0,1,0)$ و $E(0,0,1)$.

\emph{1.} للمستوي $(BDE)$ \omterm{def:g10:algebra:equation}{المعادلة} $x + y + z = 1$ (وتحققها النقط
الثلاث، وهي ليست على استقامة واحدة)، إذن $\vec n(1,1,1)$ ناظمة
له. وبما أن $\vect{AG} = (1,1,1) = \vec n$، فإن القطر $(AG)$ عمودي
على المستوي $(BDE)$.

\emph{2.} المسافة من $A(0,0,0)$:
$\frac{\abs{0 + 0 + 0 - 1}}{\sqrt3} = \frac{1}{\sqrt3} = \frac{\sqrt3}{3}$.
والمستقيم $(AG)$ هو $(t, t, t)$؛ وهو يلقى المستوي عندما يكون $3t = 1$: عند
النقطة $\left(\frac13, \frac13, \frac13\right)$، وهي مركز ثقل
المثلث $BDE$.
\end{solution}

\begin{solution}{exo:g12:space:9}
\emph{1.} $\vect{AB}(1, 0, -1)$ و $\vect{AC}(-1, 1, 0)$
و $\vect{AD}(1, 1, 1)$. فإذا كان $\vect{AD} = a\vect{AB} + b\vect{AC}$، أعطت
\omterm{def:g10:coordgeom:system}{الإحداثيات} $a - b = 1$ و $b = 1$ و $-a = 1$: فالأوليان تفرضان
$a = 2$، وهذا يناقض $a = -1$. إذن ليست \omterm{def:g12:space:collinear}{متوافقة في مستوٍ}.

\emph{2.} $\vect{BC}(-2, 1, 1)$ و $\vect{BD}(0, 1, 2)$. وبحل
$\vec n \cdot \vect{BC} = \vec n \cdot \vect{BD} = 0$: من
$b + 2c = 0$ نجد $b = -2c$؛ ثم يعطي $-2a - 2c + c = 0$ أن $c = -2a$. ومع
$a = 1$: $\vec n(1, 4, -2)$. والمستوي المار بالنقطة $B(2,1,0)$:
\[
x + 4y - 2z - 6 = 0 .
\]

\emph{3.} المسافة من $A(1,1,1)$:
$\dfrac{\abs{1 + 4 - 2 - 6}}{\sqrt{21}} = \dfrac{3}{\sqrt{21}}$.
ومساحة $BCD$: $\norm{\vect{BC}}^2 = 6$ و $\norm{\vect{BD}}^2 = 5$
و $\vect{BC}\cdot\vect{BD} = 3$، إذن
\[
\text{المساحة} = \frac12\sqrt{6 \times 5 - 9} = \frac{\sqrt{21}}{2}.
\]

\emph{4.}
\[
V = \frac13 \times \frac{\sqrt{21}}{2} \times \frac{3}{\sqrt{21}}
= \frac12 .
\]
\end{solution}

\begin{solution}{pb:g12:space:1}
\textbf{1.} المستوي: $x + y + z = 1$. والمستقيم: $(t, t, t)$.
\omterm{def:g10:numbers:interunion}{والتقاطع}: $3t = 1$:
$\left(\frac13, \frac13, \frac13\right)$.

\textbf{2.} $\dfrac{\abs{0 + 0 + 0 - 1}}{\sqrt3} =
\dfrac{1}{\sqrt3} = \dfrac{\sqrt3}{3}$.

\textbf{3.} $1 - 1 + 0 = 0$: إذن متعامدتان.

\textbf{4.} الناظمة نفسها $(1,1,1)$، لكن $\frac52 \neq 1$:
إذن مستويان متوازيان تمامًا. وأما المستقيم: فتوجيهه
$(1,1,0) \cdot (1,1,1) = 2 \neq 0$: وهو غير موازٍ للمستوي،
فيخترقه في نقطة واحدة.

\textbf{5.} المنتصفات: $\left(1, \frac12, 0\right)$ و
$\left(\frac12, 1, 0\right)$ و $\left(0, 1, \frac12\right)$ و
$\left(0, \frac12, 1\right)$ و $\left(\frac12, 0, 1\right)$ و
$\left(1, 0, \frac12\right)$: ومجموع \omterm{def:g10:coordgeom:system}{إحداثيات} كل منها
$\frac32$.

\textbf{6.} ناظمة $P$ هي $(1,1,1)$، وهي توجيه
$(AG)$: إذن متعامدان. والمركز
$\left(\frac12,\frac12,\frac12\right)$ مجموع إحداثياته
$\frac32$: فهو على $P$.

\textbf{7.} تختلف المنتصفات المتعاقبة بمتجهات مثل
$\left(-\frac12, \frac12, 0\right)$: وطولها
$\frac{\sqrt2}{2}$ في كل مرة --- أي ستة أضلاع متساوية.

\textbf{8.} عند $\left(\frac12, 1, 0\right)$: متجهتا الحافتين
$\left(\frac12, -\frac12, 0\right)$ و
$\left(-\frac12, 0, \frac12\right)$: وجداؤهما السلّمي $-\frac14$،
ومعياراهما $\frac{\sqrt2}{2}$: $\cos = -\frac12$: فالزاوية
$120^\circ$. وأضلاع متساوية وزوايا $120^\circ$ متساوية: أي مسدَّس
منتظم.

\textbf{9.} المحيط $6 \times \frac{\sqrt2}{2} = 3\sqrt2
\approx 4.24$. والمساحة $6 \times \frac{\sqrt3}{4} \times
\frac12 = \frac{3\sqrt3}{4} \approx 1.30$: أي أكبر من وجه
(مساحته $1$) --- شريحة أكبر من أي جانب من الصندوق؛ ولا
يتغلب عليها إلا المستطيل القطري ($\sqrt2 \approx 1.41$).

\textbf{10.} يلقى $(t,t,t)$ المستوي $P$ عند $t = \frac12$: أي المركز
--- وهو \omterm{prop:g10:coordgeom:midpoint}{منتصف} $[AG]$ ($A$ عند $t = 0$ و $G$ عند
$t = 1$).

\textbf{11.} من أجل $c = \frac12$: يقطع المستوي الحواف الثلاث
عند زاوية $A$، عند $\left(\frac12,0,0\right)$ و
$\left(0,\frac12,0\right)$ و $\left(0,0,\frac12\right)$: أي
مثلث متساوي الأضلاع طول ضلعه $\frac{\sqrt2}{2}$. ولما ينمو $c$،
ينمو المثلث، وتُشذَّب زواياه فتصير مسدَّسًا
(منتظمًا بالضبط عند $c = \frac32$)، ثم تنكمش الصورة
بتناظر عائدةً إلى مثلث عند زاوية $G$: وهو
تحول القطع البلوري.

\textbf{12.} كل ضلع من المقطع هو \omterm{def:g10:numbers:interunion}{تقاطع}
مستوي القطع مع \emph{وجه} واحد من المكعب، والمستوي
يلقى وجهًا (وهو مضلع مستوٍ) في قطعة واحدة على الأكثر:
أي $6$ أضلاع على الأكثر. فالسبعة مستحيلة؛ وأما من ثلاثة إلى ستة فكلها
تحدث (والسؤالان 8 و 11 يبيّنان اثنين منها).

\textbf{13.} القاعدة $ABD$: مثلث قائم مساحته $\frac12$؛
والارتفاع $AE = 1$: فالحجم
$\frac13 \times \frac12 \times 1 = \frac16$.

\textbf{14.} كل زوج من بين $B, D, E, G$ يختلف بالضبط في
إحداثيين بمقدار $1$: فالمسافات الست كلها تساوي $\sqrt2$ ---
أي منتظم. وينقسم المكعب إلى $BDEG$ زائد أربعة رباعيات وجوه
في الزوايا مطابقة للرباعي $ABDE$: فالحجم
$1 - 4 \times \frac16 = \frac13$.

\textbf{15.} المستوي $(BDE)$: $x + y + z = 1$؛ والمركز
$\left(\frac12,\frac12,\frac12\right)$: فالمسافة
$\frac{\abs{\frac32 - 1}}{\sqrt3} = \frac{1}{2\sqrt3} =
\frac{\sqrt3}{6}$.

\textbf{16.} $\vect{AG} = (1,1,1)$ و
$\vect{BH} = (-1,1,1)$: $\cos\theta = \frac{-1 + 1 + 1}{3} =
\frac13$: $\theta \approx 70.5^\circ$، ومتممتها
$\approx 109.5^\circ$. وأربعة أزواج إلكترونية يتنافر بعضها مع بعض
وتتباعد قدر الإمكان حول الكربون: فالأمثل
هو زوايا رباعي الوجوه المنتظم، الذي تلتقي اتجاهاته من
المركز إلى الرؤوس عند هذه الزاوية بالضبط ---
فالألماس هو السؤال 14 متبلورًا.

\textbf{17.} $(1, t, t)$ على $P$: $1 + 2t = \frac32$:
أي $t = \frac14$: فالنقطة $\left(1, \frac14, \frac14\right)$.

\textbf{18.} $(AB)$: نقطه $(t, 0, 0)$؛ و $(EG)$: نقطه
$(s, s, 1)$. واللقاء سيقتضي $0 = 1$ في الإحداثي
الثالث: وهذا لا يكون أبدًا. والتوازي سيقتضي أن تكون $(1,0,0)$ و
$(1,1,0)$ \omterm{def:g12:space:collinear}{مستقيميتين}: وهذا غير حاصل. فلا لقاء ولا توازٍ: إذن متخالفان
--- ممرّان في طابقين مختلفين.

\textbf{19.} $B + t(1,1,1) = (1 + t, t, t)$ على $P$:
$1 + 3t = \frac32$: أي $t = \frac16$: فالمسقط
$\left(\frac76, \frac16, \frac16\right)$، ومجموع إحداثياته
$\frac32$: فهو على $P$، كما هو مطلوب.

\textbf{20.} المستقيمات الوسيطية تخترق (السؤالان 1 و 10 و 17)؛
والمتجهات الناظمة تقيس الزوايا والمسافات (الأسئلة 2 و 8 و
15 و 16)؛ ومعادلات المستويات تنحت المقاطع (تحوّل المثلث
إلى مسدَّس). وقد ردّ الملعب زيارتنا: فمسدَّس
منتظم في صندوق من المربعات، ورباعي وجوه منتظم في
الزوايا، وزاوية الألماس، ومستقيمان يتجاهل أحدهما الآخر
--- وهي هندسة لا يقدمها إلا الفضاء.
\end{solution}
